%\chapter*{Preface to Crystallographic Powder Diffraction Analysis with GSAS-II}
\chapter{Preface to the preview version}
I am something of a perfectionist. Either that, or I am overly embarrassed to have people see my mistakes. 
Normally, I proofread anything I publish until I find myself changing what I have already changed multiple times. 
However, here I have decided to do something different, since I am releasing this as an open-source textbook and I expect there will be many updates to come. 
So, I am making this text available in a preview version (Edition 0.9), even though I have not yet started reviewing it. 
I am sure that when I do, I will find plenty I'd like to correct. 
I'm also sure that I will think of more material that I will want to add.
However, I also know that there are people who can make use of this information now, in preference to making it available in six or perhaps sixty months, when I am comfortable that this is ``good enough,'' even if not perfect.

Knowing that this can be updated allows me to be a bit more comfortable with the idea of allowing you, the reader, to see something that I consider unfinished. I do ask that you let me know of the errors that you see, be they typos, bad grammar, or wrong math, but please be gentle with my embarrassment. 
Please let me know what sections are hard to follow or perhaps you find incomprehensible. I also want to know what information is missing -- what would you like to see covered that is not here. Use \url{https://www.anl.gov/staff-directory} to send me an e-mail or contact me on GitHub. \\
\\
Brian Toby

\chapter{Preface}

\section*{Why am I writing this book?}

 As I approach the later stages of my career, I can reflect on the learning curve that I experienced in getting started with Rietveld refinement, as well as what I have experienced as an instructor and as a technique developer. 
 I interact frequently with scientists and future scientists trying to use Rietveld fitting to determine something in their research work. Periodically find myself wanting to rant, ``Rietveld cannot be used like a black box, you need to know what the parameters do and add them to the refinement slowly and only the ones that are needed.''
 The problem with that statement is that I cannot point to a single comprehensive text that will tell the recipient: how Rietveld analysis works, how our code (GSAS-II) works, what the parameters mean, and how to conduct a refinement. 
 I was lucky to have had my hand held by some of the giants of the field when I was starting with Rietveld. I was also fortunate to have a few years of experience with single crystal structure analysis, back before that became a sometimes black box technique, which gave me a very good start on understanding crystallographic analysis. 
 Perhaps someday, the Rietveld analysis codes will be good enough that they can be used as a black box, without understanding the theory and implementation behind the practice. 
 Until then, you really must understand what the software is doing, and why the different options are present, in order to obtain meaningful results. 
 My expectation is that these computer codes will continue to improve; the rate of this improvement will continue to accelerate. However, human guidance will still determine the difference between good and less than adequate scientific results.

\section*{What is crystallographic analysis about?}

 Crystals are solids where the atoms or molecules that comprise that material are arranged in regular patterns. Crystallography uses indirect information obtained by scattering quanta, most commonly x-ray photons, neutrons and electrons, from the crystal, to determine how the atoms comprising the structure of the material are arranged. Fundamentally, knowing this atomistic structure of a material is essential to understand what the substance is, be it a pharmaceutical, a catalyst, a protein or a battery or solar cell component. Knowing how the atoms are arranged in a material is an essential part of learning how that material functions. With powder diffraction, when a material that contains more than one crystalline phase, the relative abundance of these phases can be determined, which can also be a key aspect of how the material functions. There are also other more macroscopic types of information that can be obtained from this scattering information, which may also be important for how the material functions under real-world conditions. 

 

\section*{About this book}

 The goal of this book is to introduce how to perform crystallographic analysis and related techniques from powder diffraction to graduate and undergraduate students, as well as perhaps some motivated high-school students, as well as, create a resource for working scientists and engineers who wish to broaden the work they perform. This book will not attempt to present all of the fundamentals needed to perform this work, as there are already plenty of introductory texts that present concepts such as the basic physics of diffraction, symmetry, lattices and so on. The ideal audience for this will be someone who is already familiar with the theory and process of analysis of single crystal diffraction data, since many concepts carry over directly to powder diffraction, but this experience is not a requirement. While crystallographic analysis from powder diffraction is more difficult than that of single crystal work, I find that the visual feedback available from graphical display of the data and the fitted pattern to offer help in understanding the analysis process, something that is missing from single crystal analysis.

\section*{Why GSAS-II?}

Computer software is required at nearly every stage in modern crystallographic analysis, though anyone performing this analysis should understand, at least in general terms, what the software is doing so that it will be used correctly. 
This book will explain how to use the General Structure and Analysis System-II (GSAS-II) software for the analysis steps presented here. 
Since I am an author of GSAS-II, it is obviously the software I know best and I am very proud of what my co-author in that project, Robert Von Dreele\index{Von Dreele, R. B.}, and I have accomplished with it, but there are a number of other reasons for choosing this package. GSAS-II is available to everyone for free, \textit{including all source code}, and it runs on the most commonly available personal computers, on Windows, MacOS and Linux. Even on the Raspberry Pi. 
There are many fine packages that can be used for analysis of powder diffraction data, but GSAS-II is  comprehensive to all types of crystallographic analysis. Also, it is the only major crystallographic package to be initiated in the current century. Finally, it is my belief that GSAS-II offers the greatest number of features. 

GSAS-II is being actively developed, with corrections and improvements being added frequently, sometimes daily. Except for small sections of the code where speed must be enhanced, it is written in Python, a modern computer language. I think knowledge of Python is a good adjunct skill for a scientist or engineer to acquire, so should you be so motivated, you are encouraged to experiment with a copy of the GSAS-II code to learn how it works and try making changes to it. The full Python interpreter and all the source code for GSAS-II are provided. 
If this results in something you wish to share, GitHub offers the opportunity to make a ``pull request,'' where you use that platform to propose changes. These are welcome for GSAS-II. 

\section*{About this book's style}

 It is customary in science to write in the passive voice (“it is best to…”) and to avoid use of pronouns. I look at this book as a conversation between me, the author, and you, the reader, which is not enhanced by such indirect language. My hope is that avoiding the passive voice will make this easier to read. Humor is hard to convey in print, but I will attempt to allow that into our conversation, lest otherwise I produce an insomnia cure. 
 
 In the scientific literature, one should provide citations wherever possible. I am not going to do that here. It would not improve the readability of this book and I would spend a lifetime researching the literature to cite, but I will include references to materials that may help in finding more details on different topics.

 I have made the decision to release this book, like GSAS-II, as open source and make it freely available. I do this knowing that an inexpensive book by the standards where I live and work would be inaccessibly expensive in some communities. I also note that I have had the good fortune to be employed through a career that has been funded largely by US citizen's \textit{and non-citizen's} tax payments. So, this is a chance for me to give something back. Finally, in making this open source, I feel more comfortable with providing access to something I consider a work in progress. What you have in front of you is not all that what I want to present and nor is it as well edited yet as I would prefer, but I will not wait until it is ``done'' since I know it can quickly be updated. 

\section*{Feedback is requested}

 As I mentioned, I see this book as a conversation, which is a two-way exchange of information. I invite readers to inform me of anything they feel is incomplete, incomprehensible or just plain incorrect. Recommendations for useful texts and web pages to reference would be helpful too. This is best done on the book's GitHub\index{GitHub} site, \url{https://github.com/briantoby/PowderCrystallography}. Reporting of typos, spelling and grammar errors is welcome as well. For those with GitHub experience – please do consider making it easy for me and submit a pull request with suggested changes for this book's text. I’m also interested to hear what topics that are not covered that you would like to see included in a future edition. Please use the GitHub issues page or discussions page. 
